\clearpage
\appendix
\section*{Appendix}
\renewcommand{\thefigure}{A\arabic{figure}}
\renewcommand{\thetable}{A\arabic{table}}
\renewcommand{\thealgorithm}{A\arabic{algorithm}}
\setcounter{figure}{0}
\setcounter{table}{0}
\setcounter{algorithm}{0}

\section{Algorithm Description}

This section provides a pseudocode description of the proposed method.

\begin{algorithm}[htb]
\caption{Algorithm for prior knowledge generation}
\label{alg:kgen}
\textbf{Input $Prompt_{n}$,$Prompt_{a}$}: Human-written prompt text\\
\textbf{Parameter $LLM$}: Large language model\\
\textbf{Parameter $Encoder$}: Pre-trained models for text embedding\\
\textbf{Parameter $N_{n}$}: Number of normal prior knowledge items to generate\\
\textbf{Parameter $N_{a}$}: Number of anomalous prior knowledge items to generate\\
\textbf{Output}: Features of prior knowledge \\
\begin{algorithmic}[1] %[1] enables line numbers
\STATE Let $i=0$,$j=0$
\STATE Let $normalTexts$, $anomalyTexts$
\WHILE{i $ < N_{n}$}
\STATE $normalTexts$.add($LLM(Prompt_{n})$)
\STATE $i += 1$
\ENDWHILE
\WHILE{j $ < N_{a}$}
\STATE $anomalyTexts$.add($LLM(Prompt_{a})$)
\STATE $j += 1$
\ENDWHILE

\STATE \textbf{return} $Encoder(normalTexts)$ 
\STATE \quad \quad \quad $Encoder(anomalyTexts)$
\end{algorithmic}
\end{algorithm}

Algorithm \ref{alg:kgen} describes how to generate and embed prior knowledge in text form. The algorithm takes domain-specific textual prompts and generates descriptions of what is normal and what is abnormal. Because random hints are added, the LLM output has a low repetition rate, so this algorithm can generate a large amount of prior knowledge.

\begin{algorithm}[htb]
\caption{Algorithm for maintaining consistency in frame descriptions}
\label{alg:maintain}
\textbf{Input $Prompt$}: Human-written prompt text for fixing errors in description\\
\textbf{Input $Frames$}: Pending video frames\\
\textbf{Parameter $LLM$}: Large language model\\
\textbf{Parameter $BLIP$}: Pre-trained model for image captioning\\
\textbf{Parameter $YOLO$}: Pre-trained object detection model \\
\textbf{Output}: Corrected video frame description text \\
\begin{algorithmic}[1] %[1] enables line numbers
\STATE Let $i=0$,$Corrected$
\WHILE{i $ < len(Frames)$}
\STATE description $= BLIP(Frames[i])$
\STATE extra $= YOLO(Frames[i])$
\STATE $Corrected$
\STATE \quad.add($LLM([Prompt, description, extra])$)
\STATE $i += 1$
\ENDWHILE

\STATE \textbf{return} $Corrected$
\end{algorithmic}
\end{algorithm}

Algorithm \ref{alg:maintain} describes how to maintain the consistency and correctness of the textual descriptions of video frames. The effectiveness of our method relies on the assumption that information produced by object detection models is more accurate. Therefore, we can use it and an LLM to correct the descriptions generated by BLIP.

\begin{algorithm}[htb]
\caption{Algorithm for cross-modal attention with prior knowledge}
\label{alg:cross}
\textbf{Input $F_{t}$}: Text features\\
\textbf{Input $F_{v}$}: Visual features\\
\textbf{Parameter $F_{k}^{n}$,$F_{k}^{a}$}: Prior knowledge features\\
\textbf{Parameter $N$}: In the attention module, the number of convolutions in the convolution group\\
\textbf{Output}: Fused features\\
\begin{algorithmic}[1] %[1] enables line numbers
\STATE Let $i=0$,$j=0$
\STATE $Q_{v} = FFN(GAP(Conv2D(F_{v})))$
\WHILE{i $ < N$}
\STATE $K_{t} = FFN_{k}(Conv1D[i](F_{t}))$
\STATE $V_{t} = FFN_{v}(Conv1D[i](F_{t}))$
\STATE $F_{t \mid v} = Concat(F_{t \mid v}, Attention(Q_{v}, K_{t}, V_{t}))$
\STATE $i += 1$
\ENDWHILE
\STATE $Q_{t} = add(Conv1D(F_{t \mid v}), dim=1)$
\WHILE{j $ < N$}
\STATE $K_{v} = FFN_{k}(GAP(Conv2D[j](F_{v})))$
\STATE $V_{v} = FFN_{v}(GAP(Conv2D[j](F_{v})))$
\STATE $F_{v \mid t} = Concat(F_{v \mid t}, $
\STATE \quad \quad \quad $Reshape(Attention(Q_{t}, K_{v}, V_{v})))$
\STATE $j += 1$
\ENDWHILE
\STATE $F_{att} = DConv2D(Shuffle(Upsample(F_{v \mid t})))$
\STATE $rate_{n} = Sig(FFN_{n}(FFN(F_{k}^{n}) + GAP(F_{att})))$
\STATE $rate_{a} = Sig(FFN_{a}(FFN(F_{k}^{a}) + GAP(F_{att})))$
\STATE $F_{fused} = Concat(Select(F_{att}*rate_{n}), $
\STATE \quad \quad \quad \quad \quad \quad \quad \quad $Select(F_{att}*rate_{a}))$

\STATE \textbf{return} $F_{fused}$
\end{algorithmic}
\end{algorithm}

Algorithm \ref{alg:cross} illustrates the cross-modal attention mechanism and how prior knowledge optimizes the features.


\section{Examples of Prior Knowledge}

This section presents the generated priors for different scenarios. Due to space limitations, only a subset of the generated results is shown.

Table \ref{tab:sht-know} shows prior knowledge generated by GPT-3 for the ShanghaiTech dataset. The scenes in this dataset are relatively simple, consisting of campus pedestrian streets. In these scenes, normal behavior is usually walking, while other behaviors such as biking or playing are abnormal. Given these dataset characteristics, we carefully design the prompt text to encourage the LLM to generate descriptions that satisfy the requirements of prior knowledge as much as possible. We also emphasize randomness in the prompt text, which maximizes the variety of generated content.

Table \ref{tab:ucf-know} shows examples generated for several anomaly types in the UCF-Crime dataset. Thanks to GPT-3's large knowledge base and reasoning ability, and under the deliberate guidance of our prompt text, rich descriptions are generated for specific anomaly types. During prior knowledge generation, we noticed that some generated texts overlap to varying degrees. However, we did not filter these patterns because descriptions from different perspectives may increase the robustness of the prior knowledge, and dictionary learning may also mitigate the impact of excessive redundancy on performance.

For the distribution of generated prior knowledge, we adopt a balanced strategy. For ShanghaiTech, since the scenes and anomalies are relatively simple, we do not subdivide anomaly types and let the LLM generate an equal number of normal and abnormal samples. For UCF-Crime, the scenes are more complex and anomaly types are more diverse, so abnormal samples account for one third of all samples, with an equal number of samples for each anomaly type. This makes the characteristics of each anomaly type more distinct and helps the model detect anomalies.

\begin{table*}[htb]
\centering
\begin{tabularx}{\textwidth}{lXXX}
\hline
Dataset & Prompt Text & Normal Description                                    & Anomalous Description                                              \\ \hline
\multirow{6}{*}{ShanghaiTech} &
  \multirow{6}{4.5cm}{I need to detect the abnormality of a certain street in the campus. On this street, only normal walking is normal. Other situations, such as passing by bicycle, passing by fast running, etc., are abnormal. Please help me to randomly generate text descriptions of normal and abnormal situations. Note that it is randomly generated.} &
  "People leisurely walking down the street" &
  "A person riding a bicycle through the pedestrian area." \\
        &             & "Students chatting while strolling along the sidewalk."   & "Someone roller-skating at a high speed down the sidewalk."          \\
        &             & "A group of friends walking together at a moderate pace." & "A jogger sprinting through the street segment."                     \\
        &             & "Someone carrying a backpack and walking calmly."         & "Individuals playing a game of frisbee in the middle of the street." \\
        &             & "Students walking with their hands in their pockets."     & "Someone skateboarding recklessly down the sidewalk. "               \\
        &             & "A student wearing headphones, walking at a steady pace." & "A student running with a backpack on their back."                   \\ \hline
\end{tabularx}
\caption{Examples of prior knowledge generated by the LLM on the ShanghaiTech dataset.}
\label{tab:sht-know}
\end{table*}





\section{Detailed Experimental Setup}

The main text describes part of our experimental setup; due to space limitations, some details are provided here.

For data processing, we extract video frames at 16 FPS and resize them to 224$\times$224. For frames with a different aspect ratio, we use zero padding. We process these frames using I3D and embed frames from the same video into a tensor of shape $(T, N, C)$, where $T$ is the number of frames, $N=10$ indicates 10-crop data augmentation, and $C=2048$ is the feature dimension. The same scheme is used to extract optical flow features after preprocessing the video frames.

For text data, after obtaining frame descriptions with our method, we use BART to embed each frame description and obtain a tensor of shape $(T, C)$, where $T$ is the number of video frames and $C=2048$ is the feature dimension. Since video durations differ, we quantize $T$ for both visual and text features to 32. Our textual prior knowledge is also embedded using BART, resulting in a tensor of shape $(N, C)$, where $N$ is the number of prior knowledge items generated by the LLM and $C=2048$ is the embedding dimension. We then apply dictionary learning to obtain a tensor of shape $(N^{'}, C)$, where $N^{'}=256$.

For model configuration, we compress both visual and textual features by a factor of $\frac{1}{8}$, reducing the dimension from 2048 to 256 to lower computational cost. We also use dropout with a rate of 0.7. In our attention module, for the coarse-grained "Visual-Text Attention", we set the FFN output dimension to 256. In contrast, for the fine-grained "Text-Visual Attention", we set the FFN output dimension to 512. We set the number of convolutions in the convolution group to 8 based on our experiments. Finally, upsampling uses bilinear interpolation, and the last dilated convolution has a dilation coefficient of 4.

For training, we use the Adam optimizer with a learning rate of 0.001 and weight decay of 0.005, without learning rate decay. We set the batch size to 32 and train for 15,000 iterations. The final model is the one with the highest performance on the validation set. The configuration of our method is summarized in Table \ref{tab:train-cfg}.

\begin{table}[htb]
\centering
\resizebox{\columnwidth}{!}{%
\begin{tabular}{l|ll}
\hline
Type                               & Item                      & Value    \\ \hline
\multirow{5}{*}{Visual processing} & FPS                       & 16       \\
                                   & Resize                    & 224      \\
                                   & Padding                   & 0        \\
                                   & Augmentation              & 10-crop  \\
                                   & Feature dimension         & 2048     \\ \hline
\multirow{3}{*}{Text processing}   & Feature dimension         & 2048     \\
                                   & Quantize dimension        & 32       \\
                                   & Dictionary dimension      & 256      \\ \hline
\multirow{7}{*}{Model}             & Feature compression ratio & $\frac{1}{8}$      \\
                                   & Dropout ratio             & 0.7      \\
                                   & Coarse dimension          & 256      \\
                                   & Fine-tuned dimension      & 512      \\
                                   & Number of convolutions    & 8        \\
                                   & Upsampling                & bilinear \\
                                   & Dilated coefficient       & 4        \\ \hline
\multirow{5}{*}{Training}          & Optimizer                 & Adam     \\
                                   & Learning rate             & 0.001    \\
                                   & Weight decay              & 0.005    \\
                                   & Batch size                & 32       \\
                                   & Iteration                 & 15000    \\ \hline
\end{tabular}%
}
\caption{Detailed hyperparameter configuration of our method.}
\label{tab:train-cfg}
\end{table}

\section{Visualization of Experimental Results}

\begin{figure*}[htb]
\centering
\includegraphics[width=\linewidth]{images/app_cam.png}
\caption{Visualization of anomalous activation maps for some samples in the UCF-Crime dataset.}
\label{fig:app-cam}
\end{figure*}

This section presents additional visualization results and analyzes them.

Figures \ref{fig:app-know-ab} and \ref{fig:app-know-no} show additional visualizations that demonstrate the superiority of our Knowledge Traction module. In the "Burglary092" sample, the Knowledge Traction module corrects wrongly detected frames by suppressing the activation intensity of normal frames through prior knowledge. In the "Explosion021" example, some abnormal frames are clearly recognized as normal. A possible reason is that the color of the explosion scene is too uniform, leading to inaccurate recognition. After adding the Knowledge Traction module, this issue is effectively alleviated because, with prior knowledge, the text modality guides the model to understand the scene information in the video more comprehensively during training.

In contrast, some normal videos, such as "Normal\_Videos\_100" and "Normal\_Videos\_870", contain fast scene changes or abnormal-looking scenes that cause the model to fail. After introducing the Knowledge Traction module, the model still makes errors in these cases, but the magnitude of the detected anomalies is reduced.

\begin{figure}[htb]
\centering
\includegraphics[width=\linewidth]{images/app_know_abnormal.png}
\caption{Visualization of the impact of the Knowledge Traction (KT) module on model performance. For abnormal videos, it improves model robustness.}
\label{fig:app-know-ab}
\end{figure}

\begin{figure}[htb]
\centering
\includegraphics[width=\linewidth]{images/app_know_normal.png}
\caption{Visualization of the impact of the Knowledge Traction (KT) module on model performance. For normal videos, it reduces the model's sensitivity to scene changes.}
\label{fig:app-know-no}
\end{figure}

\begin{figure}[htb]
\centering
\includegraphics[width=\linewidth]{images/app_scatter.png}
\caption{2D visualization of features from the UCF-Crime dataset. Comparison between our method and the baseline MIL method.}
\label{fig:app-scatter}
\end{figure}

In Figure \ref{fig:app-scatter}, we visualize the features generated by the model on UCF-Crime by reducing their dimensionality. These features are labeled with 14 categories, comprising one normal label and 13 abnormal labels. The visualization shows that our model's features exhibit improved discrimination among categories. This improvement can be attributed to the effective fusion of visual and textual information within our features. Consequently, during training, our model gradually develops a deeper understanding of the distinctions between different anomaly types. This observation further highlights the benefits of multimodal fusion, as it enables the model to extract superior features from the data. Additionally, we visualize the activation map of a sample, as shown in Figure \ref{fig:app-cam}. The results reveal that our method concentrates more precisely on the region where the anomaly occurs. At the same time, this concentration is accompanied by less divergence than alternative approaches. Notably, our model successfully activates certain objects within the video frame even when they occupy a relatively small area. This emphasizes the effectiveness of our model in capturing meaningful anomalies.



% "Explosion021" there is an obvious situation where abnormal frames are recognized as normal frames. The possible reason for this situation is that the color of the explosion scene is too single, leading to inaccurate recognition. After adding the knowledge traction module, this situation is effectively alleviated because the text modality can guide the model to understand the scene information in the video sequence more comprehensively with the help of prior knowledge during the training process.





\begin{table*}[p]
\centering
\begin{tabularx}{\textwidth}{lXXX}
\hline
Scenes &
  Prompt Text &
  Normal Description &
  Anomalous Description \\ \hline
\multirow{3}{*}{Steal} &
  \multirow{3}{4.5cm}{I need to detect anomalies related to theft. The detection scenes are usually the interior scenes of shops or houses, which are rich and varied. Please help me randomly generate text descriptions of normal and abnormal scenarios.} &
  "A well-lit retail store with shelves neatly arranged, customers browsing products, and employees assisting them." &
  "A dimly lit retail store with broken shelves, items scattered on the floor, and a person acting suspiciously and concealing items in their bag." \\
 &
   &
  "A cozy living room with comfortable furniture, family members watching TV, and soft music playing in the background." &
  "A disheveled living room with overturned furniture, shattered glass, and signs of forced entry through a broken window." \\
 &
   &
  "A quiet bookstore with customers peacefully reading, classical music playing, and a barista serving coffee at the in-store cafe." &
  "A chaotic bookstore with books strewn all over the floor, a fire alarm blaring, and an individual stuffing books into a bag hurriedly." \\
  \hline
\multirow{3}{*}{Explosion} &
  \multirow{3}{4.5cm}{I need to detect anomalies about explosions. The detection scenarios are rich and diverse. The cause of the explosion could be various reasons. Please help me randomly generate text descriptions of normal and abnormal scenarios.} &
  "A bustling parking lot on a sunny day, cars coming and going as people run errands." &
  "A deafening explosion shatters the quietness of the parking lot, sending debris and smoke into the air." \\
 &
   &
  "A quiet residential street with children playing, neighbors chatting, and birds chirping." &
  "Chaos erupts on the residential street as a building collapses due to a gas leak, causing panic and confusion." \\
 &
   &
  "A serene gas station with vehicles filling up, attended by workers in their uniforms." &
  "An intense fire breaks out at the gas station, flames engulfing fuel pumps and sending plumes of black smoke into the sky." \\
  \hline
\multirow{3}{*}{RoadAccidents} &
  \multirow{3}{4.5cm}{I need to detect anomalies about road accidents. The detection scene is generally outdoors, which is rich and varied. The types of accidents are equally diverse. Please help me randomly generate text descriptions of normal and abnormal scenarios.} &
  "Vehicles smoothly flow along the avenue, obeying traffic signals and maintaining safe distances." &
  "A car suddenly speeds through red lights, narrowly missing pedestrians and causing chaos." \\
 &
   &
  "Cars are moving at a consistent speed in their respective lanes, maintaining proper following distances." &
  "A vehicle hydroplanes, loses control, and collides with the median barrier." \\
 &
   &
  "Neighbors jog and walk their dogs while cars drive by steadily, without any aggressive behavior." &
  "A drunk driver swerves erratically, posing a serious threat to everyone on the road." \\ \hline
\end{tabularx}
\caption{Examples of prior knowledge generated by the LLM for different scenarios in the UCF-Crime dataset.}
\label{tab:ucf-know}
\end{table*}
